\chapter{Algebraic constructions and exact certificates}
\label{ch:arithmetic}

The order-three Haagerup category belongs to a much larger algebraic
problem.  Can one construct associators for similar fusion rules at
other group orders?  How can special functions produce exact finite
categorical data?  And what does a computer certificate establish
when the desired conclusion concerns an infinite-dimensional CFT?
This chapter gives a guided answer.  Its purpose is to explain the
logic and several small exact calculations, rather than reproduce
the full contour analysis or a catalogue of number fields.

\section{The cyclic Haagerup--Izumi family}
\index{Haagerup--Izumi construction}

Replace $\Z_3$ by $G=\Z/n\Z$, with $n$ odd.  There are $2n$
simple objects $A_g$ and $R_g$, with fusion
\begin{align}
 A_gA_h&=A_{g+h},&A_gR_h&=R_{g+h},\notag\\
 R_gA_h&=R_{g-h},&R_gR_h&=A_{g-h}+\sum_{k\in G}R_k.
 \label{eq:HI-general-fusion}
\end{align}
The dimension equation is
\[
 z^2=nz+1.
\]
To avoid changing the meaning of the fixed order-three dimension
$d$ elsewhere in the book, write the two general roots as
\begin{equation}
 d_n=\frac{n+\sqrt{n^2+4}}2,\qquad
 \delta_n=\frac{n-\sqrt{n^2+4}}2=-d_n^{-1}.
 \label{eq:HI-general-roots}
\end{equation}
The positive Frobenius--Perron dimension is $d_n$.  The algebraic
construction described below naturally begins at the negative
spherical dimension $\delta_n$.

At $n=1$, the rule reduces to Fibonacci fusion.  The positive and
negative spherical categorifications are the Fibonacci and
Yang--Lee categories.  At $n=3$, the rule becomes the Haagerup
ring, but the associator still has to be identified: a construction
of the ring's categorification does not automatically select the
$\HH_3$ category used for the physical chain.

\section{The finite equations before the special function}

For a useful untwisted presentation, the nontrivial information is
encoded in an $n\times n$ matrix $\mathsf A_{g,h}$.  The group
labels are taken modulo $n$.  Set $\delta=\delta_n$ for this
section.  Among the required equations are rotation symmetry,
an axis sum, and the quadratic and cubic identities
\begin{align}
 \mathsf A_{g,h}&=\mathsf A_{-h,g-h}
 =\mathsf A_{h-g,-g},\label{eq:HI-rotation}\\
 \sum_h\mathsf A_{h,0}&=-\delta^{-1},\label{eq:HI-axis-sum}\\
 \sum_m\mathsf A_{m+g,h}\mathsf A_{h,m}
 &=\mathbf1_{g=0}-\delta^{-1}\mathbf1_{h=0},
 \label{eq:HI-quadratic}\\
 \sum_m\mathsf A_{m,g+h}\mathsf A_{g,m+k}\mathsf A_{h,m+l}
 &=\mathsf A_{g+l,k}\mathsf A_{h+k,l}
 -\delta^{-1}\mathbf1_{g=0}\mathbf1_{h=0}.
 \label{eq:HI-cubic}
\end{align}
The symbol $\mathbf1_P$ is the indicator of a condition $P$.
It is not an object of the category.

The quadratic identity expresses cancellation of inverse
recouplings.  The cubic identity compares the three-recoupling
and two-recoupling routes around a pentagon.  Its correction term
is important: it accounts for an additional invertible intermediate
channel when the relevant labels vanish.  Deleting that term would
give a superficially cleaner equation for the wrong fusion space.

The reconstruction theorem is usually stated with an additional
quartic identity.  In this setup the quartic equations follow by
contracting the quadratic and cubic identities.  Thus one can first
solve the conceptually simpler pentagon-type system and then match
the hypotheses of the reconstruction theorem.  The conclusion is
a complex spherical fusion category with the specified fusion
rules and dimension parameter.

These equations are polynomial identities over a field.  They do
not contain complex conjugates.  Consequently their algebraic
meaning is different from a unitary identity such as
$F^\dagger F=I$.

\begin{example}[The entire order-one coefficient problem]
When $n=1$, there is a single matrix entry
$\mathsf A=-\delta^{-1}$, with $\delta^2=\delta+1$.
The quadratic equation becomes
\[
 \delta^{-2}=1-\delta^{-1},
\]
which follows after dividing the dimension equation by $\delta^2$.
The cubic becomes
\[
 -\delta^{-3}=\delta^{-2}-\delta^{-1}.
\]
Multiplying by $\delta^3$ gives $-1=\delta-\delta^2$,
again exactly the dimension equation.  Both dimension roots solve
the polynomial conditions.  Positivity has not entered this
calculation.
\end{example}

\section{A factorized ansatz and its exceptional point}

Put $\tau_0=\delta-1$ and choose nonzero auxiliary coefficients
$W_g$ satisfying
\[
 W_0=-1,\qquad W_gW_{-g}=\delta\quad(g\ne0).
\]
A useful core is
\begin{equation}
 \mathsf A_{a,b}=
 \begin{cases}
 \displaystyle\frac{W_aW_{b-a}}{\tau_0 W_b},&b\ne0,\\[6pt]
 \displaystyle\frac{-1+\tau_0\mathbf1_{a=0}}{\tau_0},&b=0.
 \end{cases}
 \label{eq:HI-factorized-core}
\end{equation}
For nonzero $b$, the reciprocal relation permits the equivalent
product $W_aW_{b-a}W_{-b}/(\delta\tau_0)$.  Its three indices
sum to zero, making the cyclic symmetry of
\eqref{eq:HI-rotation} visible.  The point $a=b=0$ requires an
extra correction and cannot be recovered by applying the nonzero
reciprocal relation to $W_0$.

For example,
\[
 \mathsf A_{0,0}=\frac{\delta-2}{\delta-1},
 \qquad
 \mathsf A_{a,0}=\mathsf A_{0,a}=\mathsf A_{a,a}
 =-\frac1{\delta-1}\quad(a\ne0).
\]
The ansatz therefore imposes a particular fixed-axis condition.
It selects a family of categories; it is not a consequence of the
Haagerup--Izumi fusion ring alone.

\begin{example}[An order-three invariant distinguishes coefficient families]
Move the dimension parameter to its positive value
$d=(3+\sqrt{13})/2$.  The fixed-axis entry becomes
\[
 \mathsf A_{0,0}=\frac{d-2}{d-1}
 =\frac{\sqrt{13}-1}{\sqrt{13}+1}
 =\frac{7-\sqrt{13}}6.
\]
This equals the $\HH_2$ invariant entry from
Chapter~\ref{ch:haagerup}.  It differs from the $\HH_3$ entry
$(2-\sqrt{13})/3$ by $d/3$.  The factorized ansatz has therefore
imposed a restrictive associator condition even at the smallest
nontrivial group order. Equality of this one entry does not identify
the entire category with $\HH_2$ or supply a positive adjoint.
Its success in solving the Haagerup fusion ring does
not justify substituting its coefficient for the $\HH_3$ block
in the physical chain.
\end{example}

The auxiliary numbers $W_g$ are not automatically physical OPE
coefficients.  They are coordinates in a factorization used to
construct an associator.  A separate dictionary would be needed
to interpret one as a normalized observable of a CFT.

\section{Why two gamma periods differ by the group order}

Choose positive periods
\begin{equation}
 \omega_1=d_n+1,\qquad \omega_2=1+d_n^{-1}.
 \label{eq:HI-gamma-periods}
\end{equation}
The dimension equation gives three useful identities:
\begin{equation}
 \omega_1-\omega_2=n,\qquad
 \omega_1\omega_2=\omega_1+\omega_2,\qquad
 \frac1{\omega_1}+\frac1{\omega_2}=1.
 \label{eq:HI-period-relations}
\end{equation}
These are the reason for the choice.  A shift by the cyclic order
is the difference of the two special-function shifts.

\subsection{A convergent definition, including the subtraction}
\index{hyperbolic gamma function}

The special function can be introduced with ordinary complex analysis.
For positive periods write $\Omega=\omega_1+\omega_2$ and first
restrict to the open strip $0<\Re z<\Omega$. Define
\begin{equation}\label{eq:HI-gamma-integral}
 \gamma(z)=\exp\left\{-\int_0^\infty
 \left[
 \frac{\sinh((2z-\Omega)t)}
      {2\sinh(\omega_1t)\sinh(\omega_2t)}
 -\frac{2z-\Omega}{2\omega_1\omega_2t}
 \right]\frac{dt}{t}\right\}.
\end{equation}
This is a definition by an absolutely convergent integral, not a
prescription to subtract two divergent answers. Keep the terms
inside the same integrand.

To see the cancellation, put $\lambda=2z-\Omega$. Taylor expansion
at the integration endpoint gives
\[
 \frac{\sinh(\lambda t)}
 {2\sinh(\omega_1t)\sinh(\omega_2t)}
 =\frac{\lambda}{2\omega_1\omega_2t}
 +\frac{\lambda(\lambda^2-\omega_1^2-\omega_2^2)}
 {12\omega_1\omega_2}\,t+O(t^3).
\]
After subtraction and division by $t$, the integrand is bounded
as $t\to0$. At infinity its hyperbolic-function part decays
exponentially because
$|\Re(2z-\Omega)|<\Omega$, while the subtracted part is
$O(t^{-2})$ after division by $t$. These bounds are uniform on
compact subsets of the strip, giving holomorphy there by the
standard parameter-integral theorem. Exponentiation makes
$\gamma$ nonzero on this strip.

Two identities are immediate. At $z=\Omega/2$ the entire
integrand vanishes, so $\gamma(\Omega/2)=1$. Replacing $z$ by
$\Omega-z$ negates the integrand, giving
$\gamma(z)\gamma(\Omega-z)=1$. For real $z$ in the strip
the logarithm is real, hence $\gamma(z)>0$.

The remaining special-function work establishes the two shifts
and meromorphic continuation:
\begin{align}
 \gamma(z+\omega_1)
 &=2\sin(\pi z/\omega_2)\gamma(z),\notag\\
 \gamma(z+\omega_2)
 &=2\sin(\pi z/\omega_1)\gamma(z),\notag\\
\gamma(z)\gamma(\omega_1+\omega_2-z)&=1.
 \label{eq:HI-gamma-identities}
\end{align}
Their relation to the integral is visible already on the overlap
$0<\Re z<\omega_2$. Subtract the logarithms at $z+\omega_1$
and $z$, using the hyperbolic addition identity, to obtain
\[
 \log\gamma(z+\omega_1)-\log\gamma(z)
 =-\int_0^\infty
 \left[\frac{\cosh((2z-\omega_2)t)}{\sinh(\omega_2t)}
       -\frac1{\omega_2t}\right]\frac{dt}{t}.
\]
This is the convergent integral for $\log(2\sin(\pi z/\omega_2))$
on that strip, with its branch real on the real interval. Its
derivative is the usual cotangent integral, and the value at
$z=\omega_2/2$ fixes the constant $\log2$. Interchanging the
periods gives the other shift. The integral evaluations and
continuation are part of the analytic theorem quoted here; the
display shows why the sine appears rather than treating the
shift as an arbitrary functional equation.

The continued function has poles at
$-r\omega_1-s\omega_2$ and zeros at
$\Omega+r\omega_1+s\omega_2$, for $r,s\geq0$. At our periods
$\omega_1/\omega_2=d_n$ is irrational, so these are simple.
There is a useful normalization check that needs no numerical
sample. Write $R=\operatorname{Res}_{z=0}\gamma(z)$.
Taking the limit in both shifts gives
\[
 \gamma(\omega_1)=\frac{2\pi R}{\omega_2},\qquad
 \gamma(\omega_2)=\frac{2\pi R}{\omega_1}.
\]
Their product is one by reflection. Positivity as $z\to0^+$
chooses the positive residue, so
\[
 R=\frac{\sqrt{\omega_1\omega_2}}{2\pi},\qquad
 \gamma(\omega_1)=\sqrt{d_n},\qquad
 \gamma(\omega_2)=d_n^{-1/2}.
\]
The integral normalization has therefore fixed the normalization
needed by the finite coefficient construction.

\subsection{From the strip function to cyclic samples}

Set $F(z)=1/\gamma(\omega_2+z)$, so $F(0)=\sqrt{d_n}$.
Reflection gives
\begin{equation}
 F(z)F(n-z)=1.
 \label{eq:HI-sample-reflection}
\end{equation}
Taking the ratio of the two shift equations gives
\[
 \frac{F(z+n)}{F(z)}
 =\frac{\sin(\pi z/\omega_1)}{\sin(\pi z/\omega_2)}.
\]
For a nonzero integer $a$, the reciprocal-period relation in
\eqref{eq:HI-period-relations} therefore implies
\begin{equation}
 F(a+n)=(-1)^{a+1}F(a).
 \label{eq:HI-sample-shift}
\end{equation}
At $a=0$ both sines vanish; their ratio must be replaced by its
limit $\omega_2/\omega_1=d_n^{-1}$.  Thus the origin is
exceptional.  Many errors in finite sampling arguments come from
using a generic shift formula at precisely such a removable
singularity without taking the correct limit.

With representatives $1\leq a\leq n-1$, define
\begin{equation}
 W_0=-1,\qquad
 W_a=-(-1)^{a(a+n)/2}\frac{F(a)}{\sqrt{d_n}}.
 \label{eq:HI-W-samples}
\end{equation}
The exponent is an integer because $n$ is odd.  Reflection gives
$F(a)F(n-a)=1$, and the total sign exponent in $W_aW_{n-a}$ is
\[
 \frac{a(a+n)+(n-a)(2n-a)}2=a(a-n)+n^2.
\]
Its first term is even and its second odd.  Therefore
\begin{equation}
 W_aW_{-a}=-d_n^{-1}=\delta_n,
\end{equation}
which proves exactly the reciprocal identity needed in
\eqref{eq:HI-factorized-core}.

This is a useful example of an analytic expression producing exact
finite algebraic relations.  It does not yet prove all quadratics
or cubics: those involve sums over the group rather than only pairs
of reflected samples.

\section{Fourier transform, contour identities, and reconstruction}

On $\Z/n\Z$, finite Fourier transform turns convolution into
pointwise multiplication.  For a function $f$ put
\[
 \widehat f(k)=\sum_{a=0}^{n-1}f(a)e^{-2\pi i ka/n},
 \qquad
 f(a)=\frac1n\sum_{k=0}^{n-1}\widehat f(k)e^{2\pi i ka/n}.
\]
The quadratic coefficient equations contain translated convolution
sums.  Computing the Fourier transform of a column of
$\mathsf B=\tau_0\mathsf A$ makes their cancellation accessible to the gamma
reflection identity.  Some recoupling blocks are anti-circulant,
so it would be inaccurate to say that the transform simply
diagonalizes every $F$-matrix.  It diagonalizes the relevant cyclic
translations and organizes their products.

\subsection{The analytic column identity and its domain}

Write $u_h(a)=\mathsf B_{a,h}$, with $1\leq h<n$, and choose
the principal frequency
\[
 \theta_h(k)\equiv\frac{2\pi k}{n}+\pi(h-1)\pmod{2\pi},
 \qquad -\pi\leq\theta_h(k)\leq\pi.
\]
The supplied analytic column identity is
\begin{equation}\label{eq:HI-column-transform}
 \widehat u_h(k)=-\omega_2e^{-ih\theta_h(k)/2}
 \gamma\left(\frac{\Omega+h}{2}-\frac{\Omega\theta_h(k)}{2\pi}\right)
 \gamma\left(\frac{\Omega+h}{2}+\frac{\Omega\theta_h(k)}{2\pi}\right).
\end{equation}
The right side involves only two gamma values instead of an
$n$-term sum. The frequency representative matters for the
continuous contour integral used to prove it, even though the
original finite transform is periodic in $k$. At an endpoint
$\theta=\pm\pi$, the congruence forces $h$ even, and the two
displayed values agree.

Here is the connection with residues. The kernel
\[
 K_h(z)=\frac{\pi e^{-i\theta z}}{\sin\pi z}
          \frac{F(z)}{F(z-h)}
\]
has residue $(-1)^ae^{-i\theta a}F(a)/F(a-h)$ at an integer
$a$. The choice of $\theta$ reproduces the signs of the finite
samples. There is one exceptional sample: the actual column
entry differs from its generic interpolation at $a=h$.
Accounting for it gives
\[
 \sum_{a=0}^{n-1}\operatorname{Res}_{z=a}K_h(z)
 =(-1)^h\sqrt{d_n}F(h)
   \left(\widehat u_h(k)+\omega_2e^{-2\pi ikh/n}\right).
\]
The rectangle residue theorem converts this sum to a two-gamma
integral and the remaining crossed-pole residues. Those residues
cancel the displayed exceptional-sample correction in the final
answer. Omitting the correction would give a wrong Fourier formula.

The continuation used here has a precise domain. For fixed
$|\theta|<\pi$, allow a complex column parameter $h'$ in
\[
 \left\{\Re h'>\Omega(|\theta|/\pi-1)\right\}
 \setminus\left(\omega_1\Z_{\geq0}+\omega_2\Z_{\geq0}\right).
\]
Choose first a negative real $h_0$ in this domain and a vertical
line $h_0<\Re z<0$, which separates the two pole families.
Locally fixed contour deformations, their crossed-pole corrections,
and uniform tail estimates continue the integral through the
connected domain. The final positive integer $0<h<n$ avoids
the excluded pinch set. At this final value the integral still
decays at $|\theta|=\pi$, so endpoint limits are valid. This is
the restricted continuation argument in the expanded and formal
development; an unqualified appeal to meromorphic continuation
would omit essential work.

\subsection{A quadratic system becomes one Fourier product}

Reflection and a period shift in
\eqref{eq:HI-column-transform} give
\begin{equation}\label{eq:HI-fourier-product}
 \widehat u_h(k)\widehat u_{n-h}(-k)
       =\omega_2^2e^{-2\pi ikh/n}.
\end{equation}
One can check the signs explicitly. By swapping $h,n-h$ and
complex conjugating real columns, it suffices to use odd $h$ and
$0\leq k\leq(n-1)/2$. Put $\lambda_k=\Omega k/n$.
The four gamma arguments in the product are
\[
 x_\pm=(\Omega+h)/2\pm\lambda_k,\quad
 y_1=\Omega+(n-h)/2-\lambda_k,\quad
 y_2=(n-h)/2+\lambda_k.
\]
They obey $x_++y_1=\Omega+\omega_1$ and
$x_-+y_2=\omega_1$. Reflection and the period shifts
reduce their product to
\[
 \frac{\sin\!\left(\pi[(h-n)/(2\omega_2)+\omega_1k/n]\right)}
 {\sin\!\left(\pi[(\Omega+h)/(2\omega_1)-\omega_2k/n]\right)}
 =(-1)^{(h-n)/2+k}.
\]
Indeed the two arguments divided by $\pi$ sum to the integer
$(h-n+2)/2+k$, and the denominator's argument divided by
$\pi$ lies strictly between zero and one. The exponential prefactors contribute
the same sign times $e^{-2\pi ikh/n}$, proving
\eqref{eq:HI-fourier-product}.

Now define the unscaled quadratic convolution
\[
 C_h(g)=\sum_a\mathsf B_{g+a,h}\mathsf B_{h,a}.
\]
Rotation symmetry gives
$\mathsf B_{h,a}=u_{-h}(a-h)$. Transforming in $g$ and
substituting $b=g+a$, then $c=a-h$, yields
\begin{align*}
 \widehat C_h(k)
 &=\widehat u_h(k)\sum_a u_{-h}(a-h)e^{2\pi ika/n}\\
 &=\widehat u_h(k)e^{2\pi ikh/n}\widehat u_{-h}(-k)
 =\omega_2^2=\tau_0^2.
\end{align*}
The inverse transform of the constant function is
$\tau_0^2\mathbf1_{g=0}$. Dividing by $\tau_0^2$ proves
the quadratic equation for every $h\ne0$ simultaneously.
At $h=0$ the exceptional axis is elementary:
\[
 C_0(g)=n-2\tau_0+\tau_0^2\mathbf1_{g=0}
       =\tau_0^2(\mathbf1_{g=0}-\delta^{-1}).
\]
The last equality follows from $\delta^2=n\delta+1$ and
$\tau_0=\delta-1$. This completes all quadratic equations.
The analytic input has been isolated in one column formula;
the remaining proof is finite Fourier algebra.

\begin{example}[An exact order-three matrix without numerical samples]
Take $n=3$ and $\delta=(3-\sqrt{13})/2<0$. The sample signs
give $W_1<0$ and $W_2>0$. Put
$p=W_1^2/W_2>0$, $q=W_2^2/W_1<0$. Then
\begin{equation}\label{eq:HI-n3-matrix}
 \mathsf B=\begin{pmatrix}
 \delta-2&-1&-1\\-1&-1&p\\-1&q&-1
 \end{pmatrix},\qquad \mathsf A=\frac{\mathsf B}{\delta-1}.
\end{equation}
The reciprocal sample identity gives $pq=\delta$.
The now-proved quadratic equation at $(g,h)=(1,1)$ reads
$1-p-q=0$. Thus
\[
 p=\frac{1+\sqrt{2\sqrt{13}-5}}2,\qquad
 q=\frac{1-\sqrt{2\sqrt{13}-5}}2.
\]
For a direct check, at $h=1$ the three unscaled quadratic sums
are $2+pq$, $1-p-q$, $1-p-q$. They equal
$\tau_0^2,0,0$, since
$(\delta-1)^2=\delta+2$ at $n=3$. The $h=2$ column has
the same checks with $p,q$ exchanged, and the axis case was
computed above.

This also gives an exact special-function sample. Since
$F(1)F(2)=1$ and $p=F(1)^3/\sqrt d$,
\[
 F(1)=(p\sqrt d)^{1/3},\qquad F(2)=F(1)^{-1},
\]
with the positive real cube root. This small-order identity was
deduced from the analytic quadratic theorem; it is not an
independent proof of that theorem at every order. The matrix
still has negative categorical dimension. Its quadratics alone
do not prove the cubic reconstruction identities; the next
example supplies that extra finite argument at order three.
\end{example}

\begin{example}[The complete order-three cubic system]
\label{ex:HI-n3-cubics}
Once the preceding matrix is given algebraically, its coefficient
identities can be verified without evaluating a gamma function.
Put $a=\delta-2$, $\tau_0=\delta-1$, and $\kappa=7-\delta$.
The three relations
$\delta^2=3\delta+1$, $p+q=1$, and $pq=\delta$ imply
\[
 p^2=p-\delta,\quad q^2=q-\delta,\quad
 \tau_0^2=\delta+2,\quad a\tau_0=3,\quad
 a^2=5-\delta,\quad \tau_0^3/\delta=\kappa.
\]
Multiplying the cubic equation by $\tau_0^3$ clears every
denominator. With all indices reduced modulo three, it becomes
\begin{equation}\label{eq:HI-n3-cleared-cubic}
 \begin{split}
 C_{g,h}(k,l)&:=\sum_{m=0}^2
       \mathsf B_{m,g+h}\mathsf B_{g,m+k}\mathsf B_{h,m+l}\\
 &=\tau_0\mathsf B_{g+l,k}\mathsf B_{h+k,l}
        -\kappa\mathbf1_{g=0}\mathbf1_{h=0}.
 \end{split}
\end{equation}
The exceptional correction is present for \emph{every} $k,l$
when $g=h=0$.

Three entries explain the cancellations. At $(g,h,k,l)=(1,1,0,0)$,
the sum is $-1+p-p^2=\tau_0$. At $(1,1,1,1)$ it is
$p^3-2$, and
\[
 p^3-2=-\tau_0p-\tau_0^2
       =\tau_0(q-\delta)=\tau_0q^2,
\]
as the right side requires. At $(0,0,0,0)$ the residual is
\[
 a^3-2-\tau_0a^2+\kappa=-a^2+5-\delta=0.
\]
The last cancellation would fail if the invertible-channel
correction were omitted.

Here is a complete hand check beyond those representatives.
For each fixed $g,h$, arrange the left side into a matrix with
rows $k=0,1,2$ and columns $l=0,1,2$. Expanding its three
summands and using the displayed relations gives
\[
 \begin{aligned}
 C_{0,0}&=\begin{pmatrix}
 a^3-2&-a^2+a-1&-a^2+a-1\\
 -a^2+a-1&-a^2+a-1&3a\\
 -a^2+a-1&3a&-a^2+a-1
 \end{pmatrix},\\[3pt]
 C_{0,1}&=\begin{pmatrix}
 -3&\tau_0&-\tau_0p\\
 \tau_0&-\tau_0q&-\tau_0q\\
 -3&-\tau_0p&\tau_0
 \end{pmatrix},\\[3pt]
 C_{1,1}&=\begin{pmatrix}
 \tau_0&\tau_0&3p\\
 \tau_0&p^3-2&\tau_0\\
 3p&\tau_0&\tau_0
 \end{pmatrix},\qquad
 C_{1,2}=\begin{pmatrix}
 \tau_0&-\tau_0q&-3\\
 -3&-\tau_0q&\tau_0\\
 -\tau_0p&\tau_0&-\tau_0p
 \end{pmatrix}.
 \end{aligned}
\]
Substitution into the right side of
\eqref{eq:HI-n3-cleared-cubic} gives the same matrices.
In the exceptional block, besides the already checked corner,
the only identities needed are
\[
 -a^2+a-1=\tau_0-\kappa=2\delta-8,
 \qquad 3a=\tau_0\delta-\kappa.
\]
The other blocks use $a\tau_0=3$ and
$p^3-2=\tau_0q^2$.

Finally, exchanging $(g,k)$ with $(h,l)$ preserves the equation.
Simultaneously negating all four indices exchanges $p$ and $q$
in the matrix, and the three scalar relations are invariant under
that exchange. The four tables therefore cover the classes
$\{(0,0)\}$,
$\{(0,1),(1,0),(0,2),(2,0)\}$,
$\{(1,1),(2,2)\}$, and $\{(1,2),(2,1)\}$ of ordered pairs
$(g,h)$. Each table includes all nine $k,l$ choices, proving
all $81$ cubic equations.

This is an independent finite coefficient verification after
the explicit matrix is specified. Its identification with the
gamma samples still uses the analytic argument above, and the
all-order theorem requires more work. At the positive dimension
embedding $\delta=d$, one can separately choose
$p,q=(1\pm i\sqrt{4d-1})/2$, making $\mathsf A$ Hermitian.
This matrix condition is not by itself the construction of the
full unitary associator; that passage uses the appropriate
reconstruction theorem.
\end{example}

\subsection{The further work needed for reconstruction}

The cubic residual, the difference between the two sides of
\eqref{eq:HI-cubic}, is treated similarly.  Residue calculations
convert finite sums into integrals built from four gamma factors.
An Euler-type integral transformation and a four-term contour
identity establish the required cancellation in a range where the
contours and convergence are controlled.  A finite-degree
interpolation argument, together with index symmetries of the
factorized ansatz, then recovers the remaining coefficients.
Finally, algebraic contraction gives the quartics required by
reconstruction.

Each step answers a different question.  The contour argument proves
an analytic identity on a justified domain.  Interpolation extends
the finite set of coefficient identities.  Reconstruction turns the
coefficient solution into a category.  None of these is replaced
by finding small numerical residuals for several values of $n$.

\begin{constructioninput}[Recorded formal coefficient theorem]
For every odd positive integer $n$, the specified gamma samples give
complex coefficient data satisfying the rotation and axis equations,
all quadratics, all cubics, and all reconstruction quartics. The
repository's completed Lean development constructs the hyperbolic
gamma function from its regularized integral and proves the required
analytic identities. Its final coefficient theorem does not assume
those identities as additional hypotheses.
\end{constructioninput}

The recorded verification is an all-order proof, not a collection
of finite residual checks. Its source supplies the analytic
interface by a concrete constructor and then passes that constructor
to the general coefficient theorem. The recorded build and axiom
audit report only Lean's standard logical axioms, with no admitted
proofs or additional axioms. Preparing this textbook included reading
the final declarations and their stated scope; it did not include
an independent rebuild of the entire development.

The short analytic manuscript and this formal development should
also be distinguished. A later source review identifies missing
continuation-domain justifications and an overly broad initial
parameter range in the short manuscript. The expanded analytic and
formal developments supply the more precise domain arguments.
We therefore quote the coefficient result with its recorded formal
scope, rather than treating the short contour sketch in this chapter
as an independent proof.

\begin{theorem}[Categorical consequence quoted from the uniform construction]
For every odd order $n\geq3$, the explicit gamma-sample construction
gives a complex spherical Haagerup--Izumi category with noninvertible
categorical dimension $\delta_n$.  A simultaneous change of field
embedding gives a pseudounitary category with the same fusion rules
and dimension $d_n$.
\end{theorem}

The order-one Fibonacci and Yang--Lee cases are treated separately
by the elementary calculation above. The displayed categorical
consequence additionally uses the Evans--Gannon reconstruction
theorem and the change of dimension embedding explained next.
These categorical steps are mathematical arguments outside the
formal coefficient theorem: that theorem has no conclusion asserting
the existence of a tensor category, a VOA, or a full CFT. The finite
audits discussed below are a third kind of verification, with their
own finite scopes. None of these results is a uniform construction
of unitary CFTs or an identification of an anyon-chain limit.
The order-three unitary categories used elsewhere in the book are
known separately and have explicit positive associator presentations.

\section{Why Galois conjugation does not carry an adjoint automatically}
\index{Galois conjugation}

The automorphism of the quadratic dimension field interchanging
$\delta_n$ and $d_n$ extends to an automorphism $\sigma$ of
$\C$.  One first extends it to the algebraic closure of $\Q$,
fixes a transcendence basis, and then extends to the algebraic
closure of the resulting field.  Apply this same $\sigma$ to
every coefficient of the solution.  After adjoining the dimension
variable and clearing nonzero denominators, the reconstruction
equations are polynomial over $\Q$.  The transformed coefficients
therefore satisfy those same equations with dimension $d_n$.
This argument does not require a prior proof that all gamma samples
belong to a specified number field.

It need not preserve complex conjugation:
\[
 \sigma(\overline z)\quad\hbox{need not equal}\quad
 \overline{\sigma(z)}.
\]
Thus a statement involving adjoints and positive norms is not
transported merely because its polynomial part is.  In particular,
the following are progressively different assertions:
\begin{enumerate}
\item The fusion ring has a positive dimension character.
\item A spherical category has dimensions equal to its
Frobenius--Perron dimensions; it is pseudounitary.\index{pseudounitarity}
\item The category admits a compatible positive adjoint, making
it unitary.
\item A vacuum VOA has a positive invariant Hermitian form.
\item Its irreducible modules and their intertwining operators
have the positivity needed for a unitary tensor category and a
local full theory.
\end{enumerate}
The final two assertions concern additional CFT structures and
cannot be deduced from the first three without the relevant
construction and analytic theorems.

There is also a useful obstruction on the negative side.  A
spherical structure on these fusion rules, when it exists, is
unique: the ratio of two pivotal structures acts by scalars
$t_x$ on simples.  Because $\rho^2$ contains both $\one$ and
$\rho$, one has $t_\rho^2=1=t_\rho$, so $t_\rho=1$.
The other summands and multiplication by invertibles then force
every $t_x=1$.  A category with negative $\dim\rho$ cannot
be made positive merely by choosing a different pivotal structure.

\section{What an exact finite certificate can prove}
\index{certificate!exact finite}

The chiral construction uses finite-dimensional spaces of lattice
states, exact OPE identities, and rank statements.  A certificate
should state the finite assertion, the coefficient field, and the
argument converting the recorded data into that assertion.

For example, a matrix $M$ over a number field has rank at least
$r$ if a specified $r\times r$ minor has nonzero determinant.
After clearing denominators, one can reduce modulo a prime avoiding
those denominators.  A nonzero reduced determinant proves that the
original determinant was nonzero.  This is an exact argument,
even if a floating calculation was used earlier to discover which
minor to choose.

The upper bound is a separate matter.  It can come from exact
linear relations, a known kernel, or a representation-theoretic
dimension bound.  A nonzero $r\times r$ minor alone proves
$\operatorname{rank}M\geq r$, not equality.

\begin{example}[Why every component of a coefficient algebra can matter]
Consider $R=\Q[t]/(t^2-1)$.  This is a product of two fields,
not a field.  The element $1+t$ is nonzero at $t=1$ but vanishes
at $t=-1$.  Thus testing a determinant only at the first component
does not prove invertibility in all components.

If a coefficient presentation is known to be one number field,
a single good specialization showing nonzero determinant already
proves the determinant is a nonzero field element.  If the
presentation is an algebra with several possible embeddings or
branches and the claim concerns all of them, one must control the
appropriate components.  A convenient sufficient certificate is a
nonzero norm, obtained as the product of determinant values at all
split points modulo a good prime.  It is sufficient, not necessary:
a zero reduced norm may simply mean the prime divides a nonzero
characteristic-zero norm.
\end{example}

A second useful certificate uses positivity.  If a residual vector
$v$ lies in a space with an established positive-definite Hermitian
form and exact arithmetic gives $\langle v,v\rangle=0$, then
$v=0$.  Without positive definiteness, a nonzero isotropic vector
can have zero norm.  The metric is part of the certificate's
mathematical hypotheses.

Finite verification must then be connected to an all-weight theorem.
Checking the first several graded pieces of a VOA does not exclude
a new current generated by higher fields.  The strong-generation
argument in Chapter~\ref{ch:finite_closure} supplies the implication from
specified finite products and their provenance to closure in every
weight.  Similarly, a formal proof of a rank lemma does not amount
to a formalization of the surrounding VOA, tensor-category, or
conformal-net theory.

\section{The role of arithmetic special values}

The coefficient fields, their Galois actions, and special-function
values lead naturally to arithmetic questions.  One may ask whether
the samples can be expressed through units in explicitly described
number fields, whether class-field or Stark-type constructions
organize their conjugates, and which embeddings admit positive
categorical structures.  Such questions can make exact coefficients
more intelligible and can provide new ways to construct or certify
them.

They do not, by themselves, specify a conformal spectrum or an OPE.
In particular, a cubic number field has degree three; this use of
``cubic'' is unrelated to a three-point coefficient in a CFT.
An identity between arithmetic special values becomes relevant to
the physics only after an explicit chain of identifications links
those values to the categorical or chiral data being used.

For the order-three center, the auxiliary quadratic data involve
the prime thirteen, and their isometry class is tightly constrained.
At larger odd orders the corresponding discriminant $n^2+4$ may
be composite.  The general metric-group description of the center
does not automatically say that every associator selects the same
cyclic group and quadratic form.  Structural form, evaluation of a
selected family, categorical unitarity, and CFT realization remain
separate questions.

\section{Exercises}

\begin{exercise}\label{ex:HI-periods}
Derive every identity in \eqref{eq:HI-period-relations} from
$d_n^2=nd_n+1$.  Then derive the nonzero-integer sample shift
\eqref{eq:HI-sample-shift} and compute the exceptional limit at
the origin.

\par\smallskip\noindent{\footnotesize\normalfont
\hyperref[sol:ex:HI-periods]{Solution on page~\pageref*{sol:ex:HI-periods}.}}\par
\end{exercise}

\begin{exercise}\label{ex:HI-sign}
Verify the parity calculation following \eqref{eq:HI-W-samples}
for arbitrary odd $n$.  Explain why the chosen integer
representatives must be stated when expressions such as
$(-1)^{a(a+n)/2}$ appear.

\par\smallskip\noindent{\footnotesize\normalfont
\hyperref[sol:ex:HI-sign]{Solution on page~\pageref*{sol:ex:HI-sign}.}}\par
\end{exercise}

\begin{exercise}\label{ex:exact-rank}
Give a $2\times2$ rational matrix that is invertible but becomes
singular after rounding its entries to six decimal places.  Prove
its invertibility exactly.  Explain the asymmetry between a
nonzero determinant modulo a good prime and a zero determinant
modulo that prime.

\par\smallskip\noindent{\footnotesize\normalfont
\hyperref[sol:ex:exact-rank]{Solution on page~\pageref*{sol:ex:exact-rank}.}}\par
\end{exercise}

\begin{exercise}\label{ex:component-norm}
In $\Q[t]/(t^2-1)$ compute the product of the two evaluations
of $1+t$ and of $2+t$.  Which element is invertible?  Reduce
modulo five and compare with the exact answer.

\par\smallskip\noindent{\footnotesize\normalfont
\hyperref[sol:ex:component-norm]{Solution on page~\pageref*{sol:ex:component-norm}.}}\par
\end{exercise}

\begin{exercise}\label{ex:unitarity-hierarchy}
Explain why applying a field automorphism to polynomial pentagon
equations is legitimate, while applying it to a positivity claim
requires additional work.  Give an example of an algebraic number
whose being positive in one real embedding is not preserved in
another real embedding.

\par\smallskip\noindent{\footnotesize\normalfont
\hyperref[sol:ex:unitarity-hierarchy]{Solution on page~\pageref*{sol:ex:unitarity-hierarchy}.}}\par
\end{exercise}

\paragraph{Source note.}
Evans--Gannon supply the categorical reconstruction framework.
For the uniform coefficient theorem, the recorded formal source is
\sourcefile{lean/HaagerupIzumi/Analytic/GammaConstruction.lean}, at
revision \texttt{698b66a}; its main result is
\texttt{EGEquations\_ruijsenaars}. The recorded build scope is
described in \texttt{lean/README.md}, and the short manuscript's
analytic exposition is assessed in the September 14 source review
\sourcefile{proof_revised/formalization_review.md}. The formal result
covers every odd order, whereas the later finite audits certify
their specified finite statements. Neither should be read as
formal verification of every structure in the CFT realization.
